\documentclass[12pt,a4paper]{article}
\usepackage[top=22mm,bottom=24mm,left=27mm,right=27mm]{geometry}
\usepackage{amsmath,amssymb}
\usepackage{xcolor}
\definecolor{link}{HTML}{205973}
\usepackage[colorlinks=true,linkcolor=link,urlcolor=link]{hyperref}
\hypersetup{pdftitle={Kernel evaluation and the greedy policy}}
\setlength{\parindent}{0pt}
\setlength{\parskip}{5pt}
\raggedbottom
\newcommand{\ip}[2]{\langle #1,#2\rangle}
\newcommand{\norm}[1]{\lVert #1\rVert}
\newcommand{\R}{\mathbb R}
\newcommand{\E}{\mathbb E}
\newcommand{\N}{\mathcal N}
\newcommand{\Jhat}{\widehat J}
\newcommand{\lead}[1]{\par\smallskip\noindent\textbf{#1}\enspace}

\begin{document}

{\LARGE\bfseries Kernel evaluation and the greedy policy\par}
\medskip
A finite quadratic program determines every greedy action of the fitted
value function, even at states outside the sample. Its solution gives the
fitted value up to an additive intercept. Dropping that intercept preserves
every action comparison, including ties.

\section*{The sampled fit and its kernel program}
Let \(\mathcal S\) and \(\mathcal A\) be finite, nonempty sets of states and
actions. At \(x\), action \(a\) costs \(g_{x,a}\in\R\) and moves to \(y\) with
probability \(p(x,y,a)\). The probabilities sum to one over \(y\). Write
\(\E_{x,a}[f(X')]=\sum_{y\in\mathcal S}p(x,y,a)f(y)\) for the expectation
over the next state. The discount satisfies \(0<\alpha<1\).

The map \(\Phi\) assigns each state a feature vector in a real Hilbert
space with inner product \(\ip{\cdot}{\cdot}\) and norm \(\norm{\cdot}\).
The kernel is \(K(x,y)=\ip{\Phi(x)}{\Phi(y)}\). A weight \(\mathbf z\)
and intercept \(b\in\R\) give the fitted cost
\(\widetilde J_{\mathbf z,b}(x)=\ip{\Phi(x)}{\mathbf z}+b\).

Fix a sample \(\widehat{\mathcal S}\) of \(N\ge1\) draws from a distribution
\(\pi\) with \(\pi_x>0\). Each draw has weight \(w_x=\nu_x/\pi_x\), where
\(\nu_x>0\) is its state's relevance weight. In sums over
\(\widehat{\mathcal S}\), each occurrence of a repeated state has its own
row, slack and multipliers. An unqualified sum \(\sum_{x,a}\) ranges
over draws and actions.

\lead{The fitted value.} Fix \(\Gamma>0\) and \(\kappa\ge 1/(1-\alpha)\).
Choose \(\mathbf z,b\) and one slack \(s_x\ge0\) per draw to solve
\[
\begin{aligned}
 \text{maximize}\quad &P(\mathbf z,b,s)=
  \frac1N\sum_{x\in\widehat{\mathcal S}}w_x\ip{\Phi(x)}{\mathbf z}+b
  -\frac{\kappa}{N}\sum_{x\in\widehat{\mathcal S}}s_x
  -\frac{\Gamma}{2}\norm{\mathbf z}^2\\
 \text{subject to}\quad&
 \widetilde J_{\mathbf z,b}(x)\le g_{x,a}+
   \alpha\E_{x,a}[\widetilde J_{\mathbf z,b}(X')]+s_x
 \quad\text{for every draw and action.}
\end{aligned}
\tag{P}
\]

\lead{The kernel program.} For each sampled row \((x,a)\), define
\[
 \N_{x,a}=\{x\}\cup\{y\in\mathcal S:p(x,y,a)>0\},\qquad
 q_{x,y,a}=\mathbf 1_{\{x=y\}}-\alpha p(x,y,a).
\]
These coefficients encode the difference between the current state and
its discounted successors. Form the matrix \(Q\) and vector \(R\) by
\begin{align}
 Q(x,a,x',a')&=
   \sum_{y\in\N_{x,a}}\sum_{y'\in\N_{x',a'}}
   q_{x,y,a}\,q_{x',y',a'}\,K(y,y'),\label{eq:Q}\\
 R(x,a)&=\Gamma g_{x,a}-
   \frac1N\sum_{y\in\widehat{\mathcal S}}w_y
          \sum_{u\in\N_{x,a}}q_{x,u,a}K(y,u).\label{eq:R}
\end{align}

Minimize \(F_0(\lambda)=\tfrac12\lambda^\top Q\lambda+R^\top\lambda\)
over the feasible set \(\Lambda\) given by
\[
 \lambda_{x,a}\ge0,\qquad \sum_a\lambda_{x,a}\le\frac{\kappa}{N}
   \quad\text{for each draw }x,\qquad
 \sum_{x,a}\lambda_{x,a}=\frac1{1-\alpha}.
 \tag{D}
\]
\newpage
\section*{From a kernel solution to greedy actions}
For any multiplier array \(\lambda\) and query state \(x\in\mathcal S\),
define the kernel prediction
\begin{equation}
\begin{aligned}
 \Jhat_\lambda(x)=\frac1\Gamma\biggl[
 &\frac1N\sum_{y\in\widehat{\mathcal S}} w_y K(x,y)\\[-2pt]
 &-\sum_{y,a}\lambda_{y,a}
  \bigl(K(x,y)-\alpha\E_{y,a}[K(X',x)]\bigr)\biggr].
\end{aligned}
\label{eq:kernel}
\end{equation}
The weight \(w_y\) belongs to the draw \(y\), including when the query
\(x\) lies outside the sample. For a value function \(J\), write
\[
 \operatorname{Greedy}(J,x)=\operatorname*{arg\,min}_{a\in\mathcal A}
      \{g_{x,a}+\alpha\E_{x,a}[J(X')]\}
\]
for the set of all greedy actions at \(x\).

\textbf{Proposition.} The matrix \(Q\) is positive semidefinite, and both
optimization problems attain their optima. For every minimizer
\(\lambda^*\) of \(F_0\) on \(\Lambda\) and every maximizer
\((\mathbf z^*,b^*,s^*)\) of (P), at every state \(x\),
\begin{equation}
 \widetilde J_{\mathbf z^*,b^*}(x)=\Jhat_{\lambda^*}(x)+b^*,
 \qquad
 \operatorname{Greedy}(\widetilde J_{\mathbf z^*,b^*},x)
    =\operatorname{Greedy}(\Jhat_{\lambda^*},x).
 \label{eq:claim}
\end{equation}
Thus any action selected from the kernel prediction's minimizing set
is greedy for the fitted value function.

\lead{Proof. The quadratic program.} Define the mean feature, discounted
feature difference and multiplier-weighted sum by
\[
 m=\frac1N\sum_{x\in\widehat{\mathcal S}}w_x\Phi(x),\qquad
 d_{x,a}=\Phi(x)-\alpha\E_{x,a}[\Phi(X')],\qquad
 v(\lambda)=\sum_{x,a}\lambda_{x,a}d_{x,a}.
\]
Since \(q_{x,y,a}=0\) off \(\N_{x,a}\), we have
\(d_{x,a}=\sum_{y\in\N_{x,a}}q_{x,y,a}\Phi(y)\).
Bilinearity in \eqref{eq:Q} and \eqref{eq:R} gives
\[
 Q(x,a,x',a')=\ip{d_{x,a}}{d_{x',a'}},\qquad
 R(x,a)=\Gamma g_{x,a}-\ip{m}{d_{x,a}}.
\]
Thus \(Q\) is symmetric and
\(\lambda^\top Q\lambda=\norm{v(\lambda)}^2\ge0\) for every real array.
Expanding the square relates \(F_0\) to the normalized dual objective \(G\),
\begin{equation}
 G(\lambda)=\sum_{x,a}g_{x,a}\lambda_{x,a}+
 \frac{\norm{m-v(\lambda)}^2}{2\Gamma}
 =\frac{F_0(\lambda)+\tfrac12\norm{m}^2}{\Gamma}.
 \label{eq:normal}
\end{equation}
The objectives have the same minimizers since \(\Gamma>0\).
The row cap permits an equal split of the required mass over draws,
using any one action. Thus \(\Lambda\) is nonempty. It is closed and bounded
in finite dimensions, so \(F_0\) attains its minimum.

\newpage
\section*{Why the optimal fit has this prediction}
Fix a minimizer \(\lambda^*\). Write
\(\mathbf z_\lambda=(m-v(\lambda))/\Gamma\).
For feasible \((\mathbf z,b,s)\) and \(\lambda\), completing the square in
\(G(\lambda)-P(\mathbf z,b,s)\) gives
\begin{align}
 G(\lambda)-P(\mathbf z,b,s)
 ={}&\tfrac{\Gamma}{2}\norm{\mathbf z-\mathbf z_\lambda}^2
   +\sum_{x,a}\lambda_{x,a}
       \bigl(g_{x,a}+s_x-\ip{d_{x,a}}{\mathbf z}-(1-\alpha)b\bigr)\notag\\
 &+\sum_{x\in\widehat{\mathcal S}}
       \bigl(\tfrac{\kappa}{N}-\textstyle\sum_a\lambda_{x,a}\bigr)s_x .
 \label{eq:gap}
\end{align}
The two sums are nonnegative by the constraints in (P) and (D).
We show that some primal point attains \(G(\lambda^*)\). Then every
primal maximizer has zero gap, forcing its weight to be
\(\mathbf z_{\lambda^*}\).

\lead{Attaining the bound.}
Set \(c_{x,a}=g_{x,a}-\ip{d_{x,a}}{\mathbf z_{\lambda^*}}\).
For \(\lambda\in\Lambda\), the segment from \(\lambda^*\) to \(\lambda\)
is feasible. The right derivative of \(G\) along this segment at
\(\lambda^*\) is nonnegative, which gives
\[
 \sum_{x,a}c_{x,a}(\lambda_{x,a}-\lambda^*_{x,a})\ge0 .
\]
Thus \(\lambda^*\) minimizes the linear cost
\(\sum_{x,a}c_{x,a}\lambda_{x,a}\) over \(\Lambda\). Its finite linear
programming dual is
\[
 \max_{u\in\R,\,s_x\ge0}\quad
 \frac{u}{1-\alpha}-\frac{\kappa}{N}\sum_{x\in\widehat{\mathcal S}}s_x
 \quad\text{subject to}\quad u\le c_{x,a}+s_x\quad\text{for all }x,a .
\]
Here \(u\) prices the total-mass equality and \(s_x\) prices the cap of row
\(x\). The linear minimum is feasible and finite. By finite-dimensional
linear programming duality the displayed maximum is attained at the same
value. For such \(u,s\), choose \(b=u/(1-\alpha)\). The displayed
constraints give feasibility in (P) with weight \(\mathbf z_{\lambda^*}\).
The two sums in \eqref{eq:gap} then total
\(\sum_{x,a}c_{x,a}\lambda^*_{x,a}-b+(\kappa/N)\sum_x s_x=0\).
The square vanishes as well. This proves attainment and the asserted
weight identity for every pair of optimizers.

\lead{Evaluating the weight.} For any \(\lambda\), take the inner product
of \(\mathbf z_\lambda\) with \(\Phi(x)\). The mean contributes
\(N^{-1}\sum_y w_y K(x,y)\), where \(y\) ranges over draws. Each
\(d_{y,a}\) contributes \(K(x,y)-\alpha\E_{y,a}[K(X',x)]\), using
symmetry of the real kernel. These are exactly the terms of
\eqref{eq:kernel}. Adding the intercept gives the first equality in
\eqref{eq:claim}.

\lead{Comparing actions.} Each transition row sums to one. Consequently,
\[
 g_{x,a}+\alpha\E_{x,a}[\Jhat_{\lambda^*}(X')+b^*]
  =g_{x,a}+\alpha\E_{x,a}[\Jhat_{\lambda^*}(X')]+\alpha b^* .
\]
Every action at \(x\) gains the same \(\alpha b^*\), preserving all
weak inequalities between scores.
This proves equality of the two minimizing sets in \eqref{eq:claim}.
\hfill\(\square\)

\newpage
\section*{A worked example with a tie at the sampled state}
Take two states \(s,t\) and one draw, at \(s\). Action \(A\) always moves
to \(s\), and \(B\) always moves to \(t\). At \(s\) the costs are
\(g_{s,A}=5/4\) and \(g_{s,B}=1\). Both actions cost zero at \(t\).
Use scalar features \(\Phi(s)=1,\ \Phi(t)=2\), so \(K(x,y)=\Phi(x)\Phi(y)\).
Take \(\pi=(1/2,1/2)\), \(\nu=(3/4,1/4)\) in state order \(s,t\), and
\(\alpha=1/2,\ \Gamma=1,\ \kappa=2\).

The sampled weight is \(w_s=3/2\), giving \(m=3/2\).
The feature differences are \(d_{s,A}=1-1/2=1/2\) and
\(d_{s,B}=1-(1/2)2=0\). In action order \(A,B\),
\[
 Q=\begin{pmatrix}1/4&0\\0&0\end{pmatrix},\qquad
 R=\begin{pmatrix}1/2\\1\end{pmatrix},\qquad
 F_0(r,2-r)=\frac{(r-2)^2}{8}+\frac32,\quad 0\le r\le2 .
\]
The minimum is at \(\lambda^*=(2,0)\). The recovered weight is
\(\mathbf z^*=3/2-2(1/2)=1/2\). An optimal primal point takes
\(b^*=2\) and \(s_s=0\). Both sampled constraints hold with equality,
and \(P=G(\lambda^*)=21/8\).

The kernel formula gives \(\Jhat_{\lambda^*}(s)=1/2\) and
\(\Jhat_{\lambda^*}(t)=1\), so the fitted values are \(5/2\) and \(3\).
The table gives the one-step action costs under each prediction.
\[
\begin{array}{c@{\qquad}cc@{\qquad}cc@{\qquad}c}
 &\multicolumn{2}{c}{\text{Kernel scores}}&
  \multicolumn{2}{c}{\text{Fitted scores}}&\\
 \text{State}&A&B&A&B&\text{Greedy actions}\\ \hline
 s&3/2&3/2&5/2&5/2&\{A,B\}\\
 t&1/4&1/2&5/4&3/2&\{A\}
\end{array}
\]
Adding \(b^*=2\) raises each score by \(\alpha b^*=1\). At the sampled
state both actions remain tied. At the unsampled state \(A\) remains
strictly cheaper.

\section*{Comparison with the original formulas}
Equations (11)--(14) in Bhat, Farias, Moallemi and Zheng give the fitted
value, \(Q\), \(R\) and the kernel prediction, respectively. All four
hold with the neighborhood defined in their \S2.3. Their quadratic
objective is \(F_0+S\) for a constant \(S\) independent of \(\lambda\).
Setting \(S=0\), as suggested in \S2.5, preserves the minimizing set.

The algorithm summary
in \S2.5 describes this neighborhood using only states reachable with
positive probability. To implement the sums, include the current state
even when the action cannot return to it.

In row \((s,B)\) above, the only successor is \(t\), yet \(q_{s,s,B}=1\).
Omitting \(s\) drops \(\Phi(s)\) from \(d_{s,B}\), changing it from \(0\)
to \(-1\). Then \(Q(s,B,s,B)\) changes from \(0\) to \(1\), and \(R(s,B)\)
from \(1\) to \(5/2\).

\smallskip
{\small
\textbf{Source.} N.~Bhat, V.~F.~Farias, C.~C.~Moallemi and A.~T.~Zheng,
\href{https://moallemi.com/ciamac/papers/kernel-adp-2012.pdf\#page=7}{\emph{Non-parametric Approximate Dynamic Programming via the Kernel Method}},
August 29, 2022, \S\S2.3--2.5.\par
\textbf{Formal proof.}
\href{https://github.com/nikhilpb/kadp/blob/main/lean/KernelPolicy.lean}{Lean theorem \texttt{kernelPolicy\_of\_optima}}
and \href{https://github.com/nikhilpb/kadp/blob/main/lean/Examples/KernelPolicy.lean}{two-state example}.}
\end{document}
